\documentclass{article}
\usepackage[T1]{fontenc}
\usepackage{lmodern}
\usepackage{graphicx}
\usepackage{amsmath,amssymb}
\usepackage[a4paper,margin=2cm]{geometry}
\usepackage[absolute,overlay]{textpos}
\setlength{\TPHorizModule}{1mm}
\setlength{\TPVertModule}{1mm}
\usepackage[indonesian]{babel}
\usepackage{hyperref}
\setlength{\emergencystretch}{2em}
% unit: Halaman 58

\begin{document}

% === Halaman 58 (llm) ===

\section*{3. $m$-bit LSB}

\begin{itemize}
    \item Untuk meningkatkan ukuran pesan yang disembunyikan, maka digunakan lebih dari 1 bit LSB untuk setiap \textit{byte}.
    \item Susunan bit pada setiap \textit{byte} adalah $b_7b_6b_5b_4b_3b_2b_1b_0$. Jika diambil 2-bit LSB, maka bit yang digunakan adalah bit $b_1$ dan bit $b_0$.
    
    Contoh: $110100\underline{10} \rightarrow 2$ bit LSB terakhir dipakai untuk menyembunyikan pesan.
    
    \item \textit{Trade-off}: Semakin banyak bit LSB yang digunakan, semakin besar ukuran pesan yang dapat disembunyikan, tetapi semakin turun kualitas \textit{stego-image}.
    \item Pesan dapat disembunyikan secara sekuensial atau secara acak pada \textit{pixel-pixel} di dalam citra.
\end{itemize}

\end{document}
